\section{Proposed Methodology}
\label{sec:method}

\subsection{Multi-Stage Architecture}

Figure~\ref{fig:architecture} separates the online decision path from the
offline adaptation and verification paths. The static path extracts a fixed PE
representation and applies LightGBM. A deployment gate may accept an early
decision or request behavioral observation. Once both scores are available, a
fusion rule returns a final score. TreeSHAP serves two distinct offline roles:
it contributes signals to active-learning acquisition, and it supplies
attributions for a controlled fidelity evaluation. The same explanation method
is used, but the estimands differ: acquisition asks which sample to label next;
fidelity asks whether a ranked set of coordinates affects the fitted model.

\begin{figure*}[t]
\centering
\fbox{\begin{minipage}{0.96\textwidth}
\centering
\textbf{Online evidence path}\par\medskip
PE sample $\rightarrow$ static feature vector $x$ $\rightarrow$ LightGBM
$\rightarrow \pstatic$ $\rightarrow$ uncertainty gate
$\rightarrow$ behavioral observation $z$ $\rightarrow \pdynamic$
$\rightarrow$ validated fusion rule $\rightarrow \pfused$ and decision
\par\medskip
\hrule
\medskip
\textbf{Offline adaptation and verification}\par\medskip
Unlabeled pool $\rightarrow$ TreeSHAP $\rightarrow$ $U,H,V$
$\rightarrow$ query ranking $\rightarrow$ oracle labels $\rightarrow$ retraining
\qquad$\mid$\qquad
Fixed model $\rightarrow$ top-$|\mathrm{SHAP}|$/random masks
$\rightarrow$ FF, EC, CF and replay verification
\end{minipage}}
\caption{Conceptual multi-stage framework. The verified pilot evaluates the
active-learning and fidelity branches. The gate and PID branch are specified but
require aligned real dynamic observations for empirical evaluation.}
\label{fig:architecture}
\end{figure*}

This separation prevents three timing quantities from being conflated. Static
inference time is the cost of the initial model; dynamic observation time includes
sandbox startup and the monitoring window; and fusion time is the arithmetic
replay after both scores exist. A fast fusion function cannot eliminate the time
needed to obtain $z$. Similarly, the time required to score an active-learning
pool does not measure annotation time.

\subsection{Static LightGBM and Native TreeSHAP}

LightGBM learns an additive ensemble of decision trees
\begin{equation}
F(x)=\sum_{m=1}^{M}\eta h_m(x),
\end{equation}
where $h_m$ is a tree, $M$ is the number of boosting rounds, and $\eta$ is the
learning rate \cite{ke2017lightgbm}. The binary score is obtained from the raw
logit through the logistic link. Tree models are appropriate for heterogeneous
PE feature groups and do not require the coordinates to share a unit or scale.

For a fitted booster, native TreeSHAP returns $d$ contributions and one bias
value. We verify (\ref{eq:additivity}) for every scored pool or fidelity sample.
The largest observed error was below $10^{-14}$ in both measured pipelines. All
rankings use the $d$ signed feature contributions; the bias is excluded. This
matters because including the bias as a feature would allocate the model's base
rate to the deletion budget and distort both acquisition and fidelity analyses.

\subsection{Adaptive PID Fusion Specification}

Suppose a session provides aligned pairs $(p_{\mathrm{s},k},p_{\mathrm{d},k})$ in a meaningful
order $k$. The implementation defines the discrepancy as
\begin{equation}
e_k=p_{\mathrm{s},k}-p_{\mathrm{d},k}.
\label{eq:pid-error}
\end{equation}
The state update is discrete by observation step:
\begin{align}
I_k &= I_{k-1}+e_k,\label{eq:pid-i}\\
D_k &= e_k-e_{k-1},\label{eq:pid-d}\\
u_k &= K_p e_k+K_i I_k+K_d D_k.\label{eq:pid-u}
\end{align}
The static weight is a clipped logistic transformation,
\begin{equation}
w_{\mathrm{s},k}=\operatorname{clip}
\left(\frac{1}{1+\exp(-u_k)},0.1,0.9\right),
\label{eq:pid-weight}
\end{equation}
and the fused score is
\begin{equation}
p_{\mathrm{f},k}=w_{\mathrm{s},k}p_{\mathrm{s},k}+
(1-w_{\mathrm{s},k})p_{\mathrm{d},k}.
\label{eq:fusion}
\end{equation}

The reference gains are $(K_p,K_i,K_d)=(0.85,0.12,0.35)$, but they are not
treated as established optima. The validation grid contains proportional gains
$\{0.425,0.85,1.7\}$, integral gains $\{0,0.06,0.12\}$, and derivative gains
$\{0,0.175,0.35\}$, together with simpler P and PI subsets. Fixed weights are
searched over $\{0,0.1,\ldots,1\}$. A candidate is selected by maximum validation
AUC, then lower validation FPR, then declared grid order. The decision threshold
is independently selected on validation for the target TPR.

The controller is reset for every independent \texttt{sequence\_id}, data split, seed,
and active-learning strategy. No state flows from validation into test. The
implementation contains neither an explicit $\Delta t$ nor integral anti-windup;
clipping the output weight is not the same as bounding the accumulated state. If
every sequence contains only one independent observation, the P, I, and D terms
cannot be interpreted as distinct historical benefits. A credible experiment
must therefore document why observations form a sequence and what one step
means.

Equations~(\ref{eq:pid-error})--(\ref{eq:fusion}) are a fusion heuristic, not a
stability theorem. When $\pstatic>\pdynamic$, the update tends to increase the
static weight; when $\pdynamic>\pstatic$, it tends to favor the dynamic score.
This often emphasizes the larger positive-risk estimate, but it does not
guarantee calibration or lower FPR. The required baselines are Static-only,
Dynamic-only, a 50/50 average, validation-selected fixed weighting, P, PI, PID,
and the untuned reference PID. The full design also includes a two-factor
comparison between Random/$\xfull$ querying and Fixed/PID fusion. A
difference-in-differences analysis is required before claiming interaction
between the two components.

\subsection{Explanation-Guided Active Learning}

At round $r$, the fitted model assigns probability $p_i$ to each pool sample.
Predictive uncertainty is
\begin{equation}
U_i=1-2|p_i-0.5|,
\label{eq:uncertainty}
\end{equation}
which equals one at the decision boundary and zero at the extremes. To represent
the directional structure of a signed SHAP vector, define the total positive and
negative attribution masses
\begin{equation}
A_i^+=\sum_j\max(\phi_{ij},0),\qquad
A_i^-=-\sum_j\min(\phi_{ij},0).
\end{equation}
When $A_i^++A_i^->0$, let
\begin{equation}
q_i=\frac{A_i^+}{A_i^++A_i^-}.
\end{equation}
Directional entropy is the binary entropy
\begin{equation}
H_i=-q_i\log_2 q_i-(1-q_i)\log_2(1-q_i),
\label{eq:entropy}
\end{equation}
with $H_i=0$ for a one-sided or zero attribution vector. $H_i$ is high when
large positive and negative contributions coexist. It is not the entropy of the
class probability distribution.

The third signal is signed SHAP variance across coordinates,
\begin{equation}
V_i=\frac{1}{d}\sum_{j=1}^{d}(\phi_{ij}-\bar{\phi}_i)^2,
\qquad
\bar{\phi}_i=\frac{1}{d}\sum_j\phi_{ij}.
\label{eq:variance}
\end{equation}
This quantity measures concentration/dispersion of one explanation. It is not
variance across trees, posterior samples, or independently trained models.

Because $U$, $H$, and $V$ occupy different numeric ranges, each signal is
min--max normalized over the current pool:
\begin{equation}
\widetilde{S}_i=\begin{cases}
\dfrac{S_i-S_{\min}}{S_{\max}-S_{\min}}, & S_{\max}>S_{\min},\\
0, & \text{otherwise}.
\end{cases}
\label{eq:normalize}
\end{equation}
The acquisition function is
\begin{equation}
Q_i=\alpha\widetilde{U}_i+\beta\widetilde{H}_i+
\gamma\widetilde{V}_i.
\label{eq:query}
\end{equation}
The complete policy uses $(\alpha,\beta,\gamma)=(0.3,0.4,0.3)$. Partial
ablations use $U+H=(3/7,4/7,0)$ and $U+V=(0.5,0,0.5)$; the uncertainty baseline
uses $(1,0,0)$. Random sampling assigns an independent random score. Ties are
resolved by ascending original CSV row identifier, making selection deterministic
given the score vector and seed.

One active-learning run follows six steps: (1) fit a fresh LightGBM model on the
current labeled set; (2) select a validation threshold; (3) evaluate and save
validation/test predictions; (4) compute $Q$ only for the remaining training
pool; (5) reveal labels for the top $b$ pool rows; and (6) repeat until the fixed
budget is exhausted. Test outcomes never affect queries or stopping. SHAP is
recomputed after each model update, so the acquisition geometry follows the
current classifier rather than a frozen initial explanation.

\subsection{Controlled Fidelity Measures}

The fidelity experiment fits a separate LightGBM model on the entire training
partition and computes TreeSHAP on the fixed test partition. For each $k$, it
compares the sample-specific top-$|\phi|$ mask in (\ref{eq:mask}) with 20 random
masks of the same selected-coordinate count. The model is not retrained after
masking; the experiment measures prediction sensitivity of one fitted model.

Let $p_i$ and $p_i^{(M)}$ denote original and masked probabilities. The raw
faithfulness statistic used by the thesis is
\begin{equation}
\mathrm{FF}_{\mathrm{raw}}=\frac{1}{n}\sum_{i=1}^{n}
\frac{|p_i-p_i^{(M)}|}{p_i+\epsilon},\qquad \epsilon=10^{-12}.
\label{eq:ff}
\end{equation}
It is intentionally unclipped and may exceed one when $p_i$ is small. We
therefore also report AUC before/after masking, paired AUC drop, and mean absolute
probability change. These quantities provide more interpretable context than an
isolated FF value.

Explanation consistency (EC) compares pairwise geometry. Standardize test
features with training means and standard deviations, compute the cosine
similarity matrix $C^x$, and compute the cosine similarity matrix $C^\phi$ of
signed SHAP vectors. EC is the Pearson correlation over the strict upper
triangles:
\begin{equation}
\mathrm{EC}=\operatorname{corr}
\left(\operatorname{ut}(C^x),\operatorname{ut}(C^\phi)\right).
\label{eq:ec}
\end{equation}
For $n=400$, this produces 79,800 pairs, but the pairs share samples and are not
independent observations for a naive significance test.

Contrastive fidelity (CF) describes separation of the two label-conditioned SHAP
centroids. With $\boldsymbol{c}_0$ and $\boldsymbol{c}_1$ denoting mean SHAP
vectors for the negative and positive test classes,
\begin{equation}
\mathrm{CF}=\frac{\|\boldsymbol{c}_1-\boldsymbol{c}_0\|_2}
{\max_i\|\boldsymbol{\phi}_i\|_2+\epsilon}.
\label{eq:cf}
\end{equation}
CF is a normalized geometric separation statistic. It is neither a causal metric
nor proof of explanation correctness. FF, EC, and CF use different scales and
constructs. Consequently, no OFS is reported; choosing normalization and weights
after observing results would manufacture a composite score without an empirical
meaning.

\subsection{Computational Considerations}

LightGBM prediction is proportional to the number and depth of traversed trees.
Native TreeSHAP adds pool-scoring overhead that grows with ensemble and path
complexity. The acquisition policy computes explanations only for methods that
use $H$ or $V$; uncertainty and random baselines avoid that cost. The fidelity
test requires one original prediction, one SHAP computation, three top-k masked
batches, and $3\times20$ random batches. Timing is reported only for the scope
actually measured and is not translated into endpoint latency or analyst cost.
